% !TeX program = tectonic
\documentclass[11pt,a4paper]{article}
\usepackage{fontspec}
\usepackage[russian]{babel}
\babelfont{rm}[Language=Default]{Liberation Serif}
\babelfont[Language=Default]{sf}{Carlito}
\babelfont[Language=Default]{tt}{DejaVu Sans Mono}
\usepackage{amsmath,amssymb,amsthm}
\usepackage{geometry}
\geometry{a4paper, top=2.4cm, bottom=2.4cm, left=2.4cm, right=2.4cm}
\usepackage{booktabs,tabularx,enumitem,xcolor,tcolorbox}
\tcbuselibrary{breakable,skins}
\definecolor{linkblue}{HTML}{1F4E79}
\definecolor{statgreen}{HTML}{2F6B3A}
\definecolor{goldacc}{HTML}{8B7E5A}
\usepackage[colorlinks=true,linkcolor=linkblue,urlcolor=linkblue,unicode=true]{hyperref}
\theoremstyle{plain}
\newtheorem{theorem}{Теорема}
\newtheorem{lemma}[theorem]{Лемма}
\newtheorem{proposition}[theorem]{Предложение}
\newtheorem{corollary}[theorem]{Следствие}
\theoremstyle{definition}
\newtheorem{definition}[theorem]{Определение}
\theoremstyle{remark}
\newtheorem{remark}[theorem]{Замечание}
\newtcolorbox{statusbox}[1][]{colback=statgreen!5,colframe=statgreen!75,
  title={\textbf{Сводка доказанного:} #1},fonttitle=\small,boxrule=0.7pt,arc=2pt,breakable}
\newtcolorbox{databox}[1][]{colback=goldacc!6,colframe=goldacc!80,
  title={\textbf{Данные протокола:} #1},fonttitle=\small,boxrule=0.7pt,arc=2pt,breakable}
\newtcolorbox{verifybox}[1][]{colback=linkblue!5,colframe=linkblue!80,
  title={\textbf{Верификация:} #1},fonttitle=\small,boxrule=0.7pt,arc=2pt,breakable}
\widowpenalty=10000
\clubpenalty=10000
\setlength{\parskip}{0.35em}
\newcommand{\CC}{\mathbb{C}}
\newcommand{\RR}{\mathbb{R}}
\newcommand{\ZZ}{\mathbb{Z}}
\newcommand{\QQ}{\mathbb{Q}}
\newcommand{\Ga}{\Gamma}
\newcommand{\Bfun}{\mathrm{B}}
\newcommand{\im}{\operatorname{Im}}
\newcommand{\re}{\operatorname{Re}}
\newcommand{\lcm}{\operatorname{lcm}}
\newcommand{\SNF}{\operatorname{SNF}}
\newcommand{\Jac}{\operatorname{Jac}}
\newcommand{\rank}{\operatorname{rank}}
\newcommand{\Ind}{\operatorname{Index}}
\newcommand{\disc}{\operatorname{disc}}
\begin{document}
\begin{center}
{\LARGE\bfseries Теорема 1: соотношения Римана как тождество}\\[0.4em]
{\large симметрия и положительность период-матрицы}\\[0.6em]
{\normalsize Исаев Исхак Хамзатович}\\[0.2em]
{\small Программа <<Динамический принцип>> --- лаборатория \texttt{hodge-laboratory}}\\[0.15em]
{\small Отдельная монография теоремы · Монография 4/16}
\end{center}
\vspace{0.4em}
{\color{linkblue}\hrule height 0.8pt}
\vspace{0.8em}

\section{Постановка}

Соотношения Римана --- симметрия и положительная определённость
мнимой части период-матрицы --- на кривой доказываются
\emph{тождественно}, без всякой численной проверки: численный
контроль тождества был бы слабее доказательства. Для стенда
$N=15/30$ это означает: на матрицах $91\times182$ и $406\times812$
проверять числом нечего --- доказано.

\begin{definition}[нормированная период-матрица]
В симплектическом базисе
$\mathfrak{B}=\{a_1,\dots,a_g,b_1,\dots,b_g\}$ нормируем формы
условием $\int_{a_i}\omega_j=\delta_{ij}$ и положим
$\Omega_{ij}=\int_{b_i}\omega_j$.
\end{definition}

\begin{theorem}[соотношения Римана]\label{th:riemann}
Для период-матрицы выполнены:
(1) $\Omega^{\mathsf{T}}=\Omega$;
(2) $\im\Omega$ положительно определена: для любого
$0\ne\lambda\in\RR^g$ имеет место
$\lambda^{\mathsf{T}}(\im\Omega)\lambda>0$.
\end{theorem}

\begin{proof}
Инструмент --- билинейное тождество Римана. Для фундаментального
$4g$-угольника $F$ с рёбрами $a_i,b_i$ и замкнутых форм
$\alpha,\beta$ формула Стокса с сокращением противоположных рёбер
даёт
\[
  \iint_F \alpha\wedge\beta
  =\sum_{i=1}^{g}\Bigl(
    \int_{a_i}\alpha\int_{b_i}\beta
    -\int_{b_i}\alpha\int_{a_i}\beta\Bigr).
\]
\emph{Симметрия:} при $\alpha=\omega_j$, $\beta=\omega_k$ голоморфный
$(2,0)$-объект $\omega_j\wedge\omega_k=0$
($\Omega^2_{C_N}=0$), и тождество даёт $\Omega_{jk}-\Omega_{kj}=0$.
\emph{Положительность:} при
$\alpha=\omega=\sum\lambda_k\omega_k$,
$\beta=\overline{\omega}$ имеем
$\omega\wedge\overline{\omega}=|g|^2dz\wedge d\bar z
=-2i|g|^2dx\wedge dy$; левая часть
$=-2i\int_{C_N}|g|^2dxdy$; правая с учётом
$\int_{a_i}\omega=\lambda_i$, $\int_{b_i}\omega=(\Omega\lambda^T)_i$
равна $-2i\lambda^T(\im\Omega)\overline{\lambda}^T$ (использована
симметрия пункта 1). Сокращая $-2i$:
$\lambda^T(\im\Omega)\overline{\lambda}^T=\int_{C_N}|g|^2dxdy>0$
--- для вещественного $\lambda\ne0$ это утверждение 2.
\end{proof}

\begin{remark}[почему не нужен самосопряжённый оператор]
<<Симплектический>> относится исключительно к косой форме
пересечений $J$ на целочисленной решётке $H_1(C_N,\ZZ)$ ---
объекту топологическому. В теореме участвуют ровно два построенных
объекта: решётка циклов и интегралы форм. Соотношения Римана ---
тождество, которое наша собственная период-матрица обязана
выполнять. Числа --- лишь одноразовая санити-строка вне текста.
\end{remark}

\begin{statusbox}[итог]
Теорема доказана билинейным тождеством Римана; единственная
аналитическая посылка --- формула Стокса. Следствие для программы:
форма $J$ --- поляризация, вход модуля индексов решётки Ходжа
(теорема 3).
\end{statusbox}

\begin{verifybox}[как воспроизвести]
Тождественная природа теоремы означает: верификация ---
машинное доказательство структуры, а не численный прогон. Lean-панель
проверяет целочисленные следствия (ранги $91$, $406$); лаборатория
--- санити-строку (V7: ранг $=g$, обусловленности $8.63$, $18.17$).
\end{verifybox}

\vfill
{\color{linkblue}\hrule height 0.6pt}
\vspace{0.5em}
{\small\color{gray}
Лицензия: индивидуальная исключительная лицензия Исаева Исхака Хамзатовича.
Все права на вычисления, формулы и тексты принадлежат автору.
Верификация: \texttt{python3 laboratory.py} (пункт <<Стенды>>/<<Сертификаты>>),
Lean: \texttt{verification/lean/HodgeLaboratory.lean}.
}
\end{document}
